\subsection{Xấp xỉ ánh xạ đặc trưng \tr{kernel}}

Trong những lần trước, chúng ta đã thấy những lợi ích mà phương pháp
\tr{kernel} có thể cung cấp bằng ánh xạ ngầm và hiệu quả một bài toán học từ
không gian đầu vào $\mathcal{X}$ đến không gian đặc trưng phong phú hơn
$\mathbb{H}$. Một nhược điểm tiềm ẩnkhi sử dụng phương pháp \tr{kernel} đó là
hàm \tr{kernel} cần được tính trên tất cả cặp điểm trong tập huấn luyện. Nếu
tập huấn luyện trở nên rất lớn, thì cần đến $O(m^2)$ chi phí bộ nhớ và
$O(m^2C_K)$ chi phí tính toán, trong đó $C_K$ là chi phí của một lần tính giá
trị hàm \tr{kernel}, có thể là quá lớn. Một điều khác cần xem xét đó là chi phí
dự đoán của một mô hình đã huấn luyện. Tính toán hàm \tr{kernel} hóa $h(x) =
\sum_{i = 1}^m \alpha_i K(x_i, x) + b$ cần đến $O(m)$ chi phí bộ nhớ và
$O(mC_K)$ chi phí tính toán (tất nhiên số lượng chính xác bộ nhớ và số lần tính
phụ thuộc vào số lượng \tr{vector} hỗ trợ).

Lưu ý rằng nếu ta sử dụng các \tr{vector} đặc trưng tường minh $\mathbf{x} \in
\mathbb{R}^N$, thì bài toán SVM dạng gốc có thể được áp dụng cho quá trình huấn
luyện. Bài toán dạng gốc này chỉ tốn $O(Nm)$ chi phí bộ nhớ và quá trình tính
toán chỉ cần $O(N)$ chi phí bộ nhớ và thời gian tính toán: $h(\mathbf{x}) =
\mathbf{w}\cdot \mathbf{x} + b$. Tuy nhiên, điều này chỉ có ích nếu $N < m$,
vốn thường không xảy ra khi xem xét ánh xạ đặc trưng tường minh $\Phi(x)$ sinh
bởi hàm \tr{kernel}. Ví dụ, cho một không gian đặc trưng đầu vào có $N$ chiều,
số chiều của không gian đặc trưng \tr{kernel} của \tr{kernel} đa thức bậc $d$
là $O(n^d)$. Còn của \tr{kernel} Gaussian là vô hạn. Rõ ràng, việc sử dụng ánh
xạ đặc trưng \tr{kernel} tường minh thường là không khả thi và lại nhấn mạnh
việc sử dụng hàm \tr{kernel} để tính tích trong ngầm là then chốt.

Trong phần này, ta sẽ cho thấy sự dung hòa là khả thi bằng cách xây dựng ánh xạ
đặc trưng \tr{kernel} xấp xỉ. Đó là những ánh xạ đặc trưng với số chiều $D$ do
người dùng cụ thể , $\Psi(x) \in \mathbb{R}^D$, đảm bảo $\Psi(x) \cdot \Psi(x')
\approx K(x, x')$ khi số chiều $D$ đủ lớn. Ta bắt đầu với một kết quả kinh điển
trong giải tích điều hòa.

\begin{thm}[Định lý Bochner]
  \label{thm:bochner}
  Một \tr{kernel} liên tục có dạng $K(x, x') = G(x - x')$ được định nghĩa trên
  một tập $\mathcal{X}$ \tr{compact} địa phương là xác định dương khi và chỉ
  khi $G$ là biến đổi Fourier của độ đo không âm. Tức là, \begin{equation}
    G(x) = \int_{\mathcal{X}} p(\omega)e^{i\omega\cdot x}\, d\omega,
  \end{equation}với $p$ là một độ đo không âm.
\end{thm}

\tr{Kernel} có dạng $K(x, x') = G(x - x')$ được gọi là \tr{kernel} \emph{bất
biến tịnh tiến} (shift-invariant). Nếu $G$ được chuẩn hóa tỉ lệ sao cho $G(0) =
1$, khi đó $p$ thực tế là một phân phối xác suất. Một số ví dụ về các
\tr{kernel} như thế được trình bày ở \tablename~\ref{tab:s-i-ker}. Mệnh đề tiếp
theo cung cấp biểu thức đơn giản hóa trong trường hợp \tr{kernel} nhận giá trị
thực.

\begin{table}[htbp]
  \centering
  \caption{Ví dụ về một số \tr{kernel} bất biến tịnh tiến được chuẩn hóa (định
  nghĩa trên $\mathbf{x}, \mathbf{x}^\prime \in \mathbb{R}^N$) và hàm mật độ
  tương ứng (định nghĩa trên $\boldsymbol{\omega} \in \mathbb{R}^N$).}
  \label{tab:s-i-ker}
  \import{\tablespath}{shift-invariant-kernel}
\end{table}

\begin{prop}\label{prop:expect-real-kernel}
  Cho $K$ là một \tr{kernel} giá trị nhận giá trị thực, bất biến tịnh tiến và
  $p$ là độ đo không âm tương ứng trong Định lý~\ref{thm:bochner}. Ngoài ra,
  giả sử rằng với mọi $x \in \mathcal{X}$, ta có $K(x, x) = 1$, để $p$ là một
  phân phối xác suất. Khi đó, ta có đẳng thức sau: \begin{equation}
    \expect_{\omega \sim p}\Bigl[ \bigl[ \cos(\omega \cdot x),
    \sin(\omega \cdot x) \bigr ]^\intercal \bigl[ \cos(\omega
    \cdot x'), \sin(\omega \cdot x') \bigr ] \Bigr ] = K(x, x').
  \end{equation}
\end{prop}

\begin{proof}
  Trước hết, vì $K$ và $p$ đều nhận giá trị thực nên ta chỉ dùng phần thực của
  $e^{ix}$. Tức là sử dụng $\real[e^{ix}] = \real\bigl[ \cos(x) +
  i\sin(x) \bigr] = \cos(x)$. Khi đó, \begin{equation}
    K(x, x') = \real\bigl[ K(x, x') \bigr ] = \int_\mathcal{X} p(\omega)
    \cos \bigl( \omega \cdot (x - x') \bigr )\, d\omega.
  \end{equation} Tiếp theo, sử dụng đẳng thức $\cos(a - b) = \cos(a)\cos(b) -
  \sin(a)\sin(b)$, ta có \begin{align}
    \int_\mathcal{X} p(\omega) \cos\bigl( \omega \cdot ( x - x' ) \bigr ) \,
      d\omega &= \int_\mathcal{X} p(\omega) \bigl( \cos(\omega \cdot x)
      \cos(\omega \cdot x') + \sin(\omega \cdot x) \sin(\omega \cdot x') \bigr
      ) \, d\omega \\
    &= \expect_{\omega \sim p} \Bigl[ \bigl[ \cos(\omega \cdot x), \sin(\omega
    \cdot x) \bigr ]^\intercal \bigl[ \cos(\omega \cdot x'), \sin(\omega \cdot
    x') \bigr] \Bigr ],
  \end{align} kết thúc chứng minh mệnh đề.
\end{proof}

Mệnh đề trên cung cấp một phương pháp rất đơn giản để sinh một ánh xạ
\tr{kernel} xấp xỉ $\Psi \in \mathbb{R}^{2D}$ với mỗi $D \ge 1$, được định
nghĩa với mọi $x \in \mathcal{X}$ bởi \begin{equation}
  \Psi(x) = \sqrt{\frac{1}{D}} \bigl[ \cos(\omega_1 \cdot x), \sin(\omega_1
  \cdot x), \dots, \cos(\omega_D \cdot x), \sin(\omega_D \cdot x) \bigr
  ]^\intercal,
\end{equation} với $w_i$, $i = 1, \dots, D$, là các mẫu độc lập cùng phân phối
(i.i.d.) theo độ đo $p$ trên $\mathcal{X}$ tương ứng với \tr{kernel} $K$ đang
xét. Do đó, \begin{equation}
  \Psi(x) \Psi(x') = \frac{1}{D} \sum_{i = 1}^D \bigl[ \cos(\omega_i \cdot x),
  \sin(\omega_i \cdot x) \bigr]^\intercal \bigl[ \cos(\omega_i \cdot x'),
  \sin(\omega_i \cdot x') \bigr]
\end{equation} là phiên bản thực nghiệm của kì vọng được tính trong Mệnh
đề~\ref{prop:expect-real-kernel}.  Bổ đề sau đây cho thấy ước lượng thực nghiệm
hội tụ đều trên tất cả các điểm trong tập xác định \tr{compact} khi $D$ tăng.

\begin{lem}
  \label{lem:prob-sup-bound}
  Cho $K$ là một hàm \tr{kernel} khả vi liên tục thỏa mãn các điều kiện của
  Mệnh đề~\ref{prop:expect-real-kernel} và có độ đo $p$ tương ứng. Ngoài ra,
  giả sử $\mathcal{X}$ là \tr{compact} và có số chiều $N$, $R$ là bán kính quả
  cầu Euclid chứa $\mathcal{X}$ và $\sigma_p^2 = \expect_{\omega \sim p }
  \bigl[ \lVert \omega \rVert^2 \bigr ] < \infty$. Khi đó, với $\Psi$ được định
  nghĩa như trong chứng minh Mệnh đề~\ref{prop:expect-real-kernel}, với mọi $0<
  r \le 2R$ và $\epsilon > 0$, bất đẳng thức sau đúng: \begin{equation}
    \prob \biggl[ \sup_{x, x' \in \mathcal{X}} \bigl\lvert \Psi(x) \cdot
    \Psi(x') - K(x, x') \bigr\rvert \ge \epsilon \biggr ] \le 2\mathcal{N}(2R,
    r) \exp\biggl( -\frac{D\epsilon^2}{8} \biggr ) +
    \frac{4r\sigma_p}{\epsilon}.
  \end{equation}Trong đó xác suất được lấy khi $\omega \sim p$ và
  $\mathcal{N}(R, r)$ là số lượng nhỏ nhất quả cầu bán kính $r$ cần để phủ quả
  cầu bán kính $R$.
\end{lem}

\begin{proof}
  Gọi $\mathcal{Z} = \{z\colon z = x - x',\, x, x' \in \mathcal{X} \}$, khi đó
  $\mathcal{Z}$ nằm trong quả cầu bán kính $2R$. Vì $\mathcal{X}$ đóng nên
  $\mathcal{Z}$ đóng và do đó $\mathcal{Z}$ \tr{compact}. Để cho tiện, kí hiệu
  $B = \mathcal{N}(2R, r)$ là số quả cầu bán kính $r$ nhỏ nhất phủ quả cầu bán
  kính $2R$ hay là phủ $\mathcal{Z}$ và gọi $z_j$, với $j \in [B]$, là tâm của
  các quả cầu này. Khi đó, với mọi $z\in \mathcal{Z}$, tồn tại $j$ sao cho $z =
  z_j + \delta$ với $\lvert \delta \rvert < r$.

  Tiếp theo, định nghĩa $S(z) = \Psi(x) \cdot \Psi(x') - K(x, x')$, với $z = x
  - x'$. Bởi vì $S$ khả vi liên tục trên tập \tr{compact} $\mathcal{Z}$ nên nó
  $L$-Lipschitz với $L = \sup_{z \in \mathcal{Z}} \bigl\lVert \nabla S(z)
  \bigr\rVert$.  Nếu $L < \epsilon / (2r)$ và với mọi $j \in [B]$ ta có
  $\bigl\lvert S(z_j) \bigr\rvert < \epsilon / 2$ thì bất đẳng thức sau đúng
  với mọi $z = z_j + \delta \in \mathcal{Z}$: \begin{equation}
    \label{eq:bound-easy-case}
    \bigl\lvert S(z) \bigr\rvert = \bigl\lvert S(z_j + \delta) \bigr\rvert \le
    L \lvert z_j - (z_j + \delta) \rvert + \bigl\lvert S(z_j) \bigr\rvert \le
    rL + \frac{\epsilon}{2} < \epsilon.
  \end{equation} Phần còn lại của chứng minh đó là chặn xác suất biến cố $L >
  \epsilon / (2r)$ và $\bigl\lvert S(z_j) \bigr\rvert > \epsilon / 2$.  Các xác
  suất và kì vọng ở phần sau đều tương ứng với các biến ngẫu nhiên
  $\omega_1, \dots, \omega_D$.

  Để chặn xác suất sự biến cố đầu, ta dùng Mệnh
  đề~\ref{prop:expect-real-kernel} và tính tuyến tính của kì vọng để có
  $\expect\Bigl[ \nabla\bigl( \Psi(x)
  \cdot \Psi(x') \bigr ) \Bigr ] = \nabla K(x, x')$. Khi đó, \begin{align}
    \expect[L^2] &= \expect \biggl[ \sup_{z \in \mathcal{Z}} \bigl\lVert \nabla
    S(z) \bigr\rVert^2 \biggr] \\
    &= \expect \biggl[ \sup_{x, x' \in \mathcal{X}} \Bigl\lVert \nabla\bigl(
    \Psi(x) \cdot \Psi(x') \bigr ) - \nabla K(x, x') \Bigr\rVert^2 \biggr ] \\
    &\le 2\expect \biggl[ \sup_{x, x' \in \mathcal{X}} \Bigl\lVert \nabla
    \bigl( \Psi(x) \cdot \Psi(x') \bigr ) \Bigl\rVert \biggr ] + 2\sup_{x, x'
    \in \mathcal{X}} \bigl\lVert \nabla K(x, x') \bigr\rVert^2 \\
    &\le 2\expect \biggl[ \sup_{x, x' \in \mathcal{X}} \Bigl\lVert \nabla
    \bigl( \Psi(x) \cdot \Psi(x') \bigr ) \Bigl\rVert \biggr ] + 2\sup_{x, x'
    \in \mathcal{X}} \biggl\lVert \expect\Bigl[ \nabla\bigl( \Psi(x) \cdot
    \Psi(x') \bigr ) \Bigr ] \biggr\rVert \\
    &\le 4\expect \biggl[ \sup_{x, x' \in \mathcal{X}} \Bigl\lVert \nabla
    \bigl( \Psi(x) \cdot \Psi(x') \bigr ) \Bigl\rVert \biggr ],
  \end{align} trong đó bất đẳng thức đầu (ở dòng 2 xuống dòng 3) sử dụng
  $\lVert a + b \rVert^2 \le 2\lVert a\rVert^2 + 2 \lVert b \rVert^2$ (suy ra
  từ bất đẳng thức Jensen) và tính cộng tính dưới của chặn trên nhỏ nhất. Bất
  đẳng thức thứ hai (ở dòng kế cuối xuống dòng cuối) cũng từ bất đẳng thức
  Jensen (dùng hai lần) và tính cộng tính dưới của chặn trên nhỏ nhất. Mặt
  khác, ta có thể rút gọn $\nabla\bigl( \Psi(x) \cdot \Psi(x') \bigr )$ theo $z
  = x - x'$ như sau: \begin{align}
    \nabla\bigl( \Psi(x) \cdot \Psi(x') \bigr ) &= \nabla \biggl( \frac{1}{D}
    \sum_{i=1}^D \cos(\omega_i \cdot x) \cos(\omega_i \cdot x') + \sin(\omega_i
    \cdot x) \sin(\omega_i \cdot x') \biggr ) \\
    &= \nabla \biggl( \frac{1}{D} \sum_{i = 1}^D \cos\bigl( \omega_i \cdot (x -
    x') \bigr ) \biggr ) \\
    &= -\frac{1}{D} \sum_{i = 1}^D \omega_i \sin\bigl( \omega_i \cdot (x - x')
    \bigr).
  \end{align} Kết hợp hai kết quả trước ta được \begin{align}
    \expect[L^2] &\le 4\expect \Biggl[ \sup_{x, x' \in \mathcal{X}}
    \biggl\lVert \frac{1}{D} \sum_{i = 1}^D -\omega_i \sin\bigl( \omega \cdot
    (x - x') \bigr ) \biggr\rVert^2 \Biggr ] \\
    &\le 4 \expect_{\omega_1,\dots,\omega_D} \Biggl[ \biggl( \frac{1}{D}
    \sum_{i = 1}^D \lVert \omega_i \rVert \biggr )^2 \Biggr ] \\
    &\le 4 \expect_{\omega_1,\dots,\omega_D} \biggl[ \frac{1}{D} \sum_{i = 1}^D
    \lVert \omega_i \rVert^2 \biggr ] \\
    &\le 4 \expect_{\omega}\bigl[ \lVert x \rVert^2 \bigr ] = 4\sigma_p^2,
  \end{align}suy ra từ bất đẳng thức tam giác, $\bigl\lvert \sin(\cdot)
  \bigr\rvert \le 1$, bất đẳng thức Jensen và $w_i$ là i.i.d.\ trong biến đổi
  cuối. Do đó ta có thể chặn trên cho biến cố đầu bằng bất đẳng thức Markov:
  \begin{equation}\label{eq:1st-event-bound}
    \prob\biggl[ L \ge \frac{\epsilon}{2r} \biggr ] \le \biggl(
    \frac{4r\sigma_p}{\epsilon} \biggr )^2.
  \end{equation} Để chặn cho biến cố sau, để ý rằng, theo định nghĩa, $S(z)$ là
  tổng $D$ biến i.i.d., mỗi biến có giá trị tuyệt đối chặn bởi $2/D$ (vì với
  mọi $x$ và $x'$, ta có $\bigl\lvert K(x, x') \bigr\rvert \le 1$ và
  $\bigl\lvert \Psi(x) \cdot \Psi(x') \bigr\rvert \le 1$), và $\expect\bigl[
    S(z) \bigr ] = 0$. Do đó, theo bất đẳng thức Hoeffding và chặn hợp xác suất
  các biến cố, ta có: \begin{equation}\label{eq:2nd-event-bound}
    \prob\biggl[ \exists j \in [B] \colon \bigl\lvert S(z_j) \bigr\rvert \ge
    \frac{\epsilon}{2} \bigr\rvert \biggr ] \le \sum_{i=1}^B \prob\biggl[
      \bigl\lvert S(z_j) \bigr\rvert \frac{\epsilon}{2} \biggr] \le
    2B\exp\biggl( - \frac{D\epsilon^2}{8} \biggr ).
  \end{equation} Cuối cùng theo \eqref{eq:bound-easy-case},
  \eqref{eq:1st-event-bound} và~\eqref{eq:2nd-event-bound}, và định nghĩa $B$
  ta có \begin{equation}
    \prob\biggl[ \sup_{z\in \mathcal{Z}} \bigl\lvert S(z) \bigr\lvert \ge
    \epsilon \biggr ] \le 2\mathcal{N}(2R, r) \exp\biggl(-
    \frac{D\epsilon^2}{8} \biggr) + \biggl( \frac{4r\sigma_p^2}{\epsilon}
    \biggr ),
  \end{equation} điều phải chứng minh.
\end{proof}

Một yếu tố then chốt trong biên chặn của bổ đề này là số bao phủ
$\mathcal{N}(2R, r)$, đại lượng phụ thuộc rất lớn vào số chiều $N$ của không
gian. Trong bổ đề tiếp theo dưới đây, ta sẽ làm rõ sự phụ thuộc này một
cách tường minh đối với một trường hợp đặc biệt đơn giản, mặc dù các lập luận
tương tự vẫn sẽ đúng cho các trường hợp tổng quát hơn.

\begin{lem}
  \label{lem:n-ball-bound}
  Cho $\mathcal{X} \in \mathbb{R}^N$ là \tr{compact} và $R$ là bán kính hình
  cầu nhỏ nhất chứa $\mathcal{X}$. Khi đó, ta có bất đẳng thức sau:
  \begin{equation}
    \mathcal{N}(R, r) \le \biggl( \frac{3R}{r} \biggr )^N.
  \end{equation}
\end{lem}

\begin{proof}
  Bằng cách sử dụng thể tích của quả cầu trong $\mathbb{R}^N$ ta có ngay $R^N /
  (r / 3)^N = (3R / r)^N$ là chặn trên của số lượng quả cầu bán kính $r/3$ có
  thể nằm trong quả cầu bán kính $R$ mà không cắt nhau. Xét một trường hợp số
  lượng quả cầu nhiều nhất này. Khi đó, mỗi điểm trong quả cầu bán kính $R$ này
  phải có khoảng cách tới tâm một trong các quả cầu nó chứa không quá $r$. Nếu
  không, ta có thể thêm vào một quả cầu bán kính $r/3$ khác vào quả cầu bán
  kính $R$ này mà vẫn thỏa điều kiện, mâu thuẫn với điều giả sử. \elaborate\
  Thật vậy, gọi $B$ là số quả cầu bán kính $r/3$ này. Giả sử các tâm của các
  quả cầu bán kính $r/3$ này là $c_i$ với $i \in [B]$; và điểm đang xét là $x$.
  Khi đó, ta có $\lVert x - c_i \rVert > r$ với mọi $i \in [B]$. Gọi điểm $c_0$
  sao cho trên \tr{vector} $\overrightarrow{xc_0}$ sao cho $\lVert x - c_0
  \rVert = r/3$.  Vẽ quả cầu tâm $c_0$ bán kính $r/3$. Khi đó quả cầu này không
  cắt bất kì các quả cầu bán kính $r/3$ khác. Nếu không, tồn tại điểm $y \in
  \mathbb{R}^N$ và $j \in [B]$ sao cho $\lVert c_0 - y \rVert, \lVert c_j - y
  \rVert < r/3$. Khi đó $\lVert x - c_j \rVert \le \lVert x - c_0 \rVert +
  \lVert c_0 - y \rVert + \lVert y - c_j \rVert < r$ (mâu thuẫn). Do đó khi
  thay các quả cầu bán kính  $r/3$ này bởi các quả cầu khác cùng tâm nhưng có
  bán kính $r$ thì các quả cầu này sẽ phủ quả cầu có bán kính $R$.
\end{proof}

Cuối cùng, kết hợp hai bổ đề trên, ta có thể chặn xấp xỉ mẫu hữu hạn một cách
tường minh.

\begin{thm}
  Cho $K$ là một hàm \tr{kernel} khả vi liên tục thỏa các điều kiện của Mệnh
  đề~\ref{prop:expect-real-kernel} và có độ đo tương ứng $p$. Ngoài ra giả sử
  $\sigma_p^2 = \expect_{\omega \sim p}\bigl[ \lVert \omega \rVert^2 \bigr ] <
  \infty$ và $\mathcal{X} \subset \mathcal{R}^N$. Gọi $R$ là bán kính của quả
  cầu Euclid chứa $\mathcal{X}$. Khi đó, với $\Psi \in \mathbb{R}^D$ được định
  nghĩa trong chứng minh của Mệnh đề~\ref{prop:expect-real-kernel} và mọi $0 <
  \epsilon \le 32R\sigma_p$ ta có bất đẳng thức sau: \begin{equation}
    \prob \biggl[ \sup_{x, x' \in \mathcal{X}} \bigl\lvert \Psi(x) \cdot
    \Psi(x') - K(x, x') \bigr\rvert \ge \epsilon \biggr ] \le \biggl(
    \frac{48R\sigma_p}{\epsilon}\biggr )^2\exp\biggl( - \frac{D\epsilon^2}{4(N
    = 2)} \biggr ).
  \end{equation}
\end{thm}

\begin{proof}
  Sử dụng Bổ đề~\ref{lem:prob-sup-bound} và~\ref{lem:n-ball-bound} với giá trị
  r được chọn: \begin{equation}
    r = \Biggl[ \frac{2(6R)^N \exp(-D\epsilon^2 / 8)}{ \bigl((4\sigma_p) /
    \epsilon \bigr)^2 } \Biggr ]^{2/(N + 2)},
  \end{equation}ta được biểu thức sau: \begin{equation}
    \prob\biggl[ \sup_{z \in \mathcal{Z}} \bigl\lvert S(z) \bigr\rVert \biggr ]
    \le 4\biggl( \frac{ 24R\sigma_p } { \epsilon } \biggr ) ^{(2N) / (N + 2)}
    \exp\biggl( - \frac{D\epsilon^2} { 4(N + 2) } \biggr ).
  \end{equation} Vì $32R\sigma_p / \epsilon \le 1$, số mũ $(2N) / (N + 2)$ có
  thể thay bằng 2, điều này kéo theo điều phải chứng minh.
\end{proof}

Định lý trước đó đưa ra một sự đảm bảo rằng có thể tìm được một ước lượng tốt
cho hàm Kernel với xác suất cao, bằng cách lấy mẫu một số lượng hữu hạn các tọa
độ $D$. Cụ thể, để đạt được một sai số tuyệt đối tối đa là $\epsilon$, ta chỉ
cần lấy mẫu một số lượng tọa độ là $D = O\bigl( N/\epsilon^2 \cdot \log(
R\sigma_p / \epsilon ) \bigr)$.
